\chapter{Brownian Motion}

\begin{center}\itshape
Independent normal numbers at the dyadic points, with shrinking variances,\\
produce a uniform limit. The limit is a Brownian motion.
\end{center}

\section{Definition and construction}

\begin{definition}[Brownian motion]\label{def:bm}
A Brownian motion satisfies $B_0=0$, has continuous paths, has independent increments, and gives $B_t-B_s$ the law $N(0,t-s)$.
\end{definition}

\begin{proposition}[covariance]\label{prop:bmgauss}
For a continuous process with $B_0=0$, independent normal increments of variance $t-s$ are equivalent to centred Gaussian finite-dimensional distributions with $\operatorname{Cov}(B_s,B_t)=\min(s,t)$.
\end{proposition}
\begin{proof}
Increments of that covariance are uncorrelated and jointly normal, hence independent by Theorem~\ref{thm:gaussvec}. The converse adds the variances of the increments up to the earlier time.
\end{proof}

Let $Z_t$ be the independent standard normals of Proposition~\ref{prop:manyuniforms}, indexed by the dyadic points of $[0,1]$. Set $B_0=0$, $B_1=Z_1$, and at each new midpoint of level $n$,
\[
 B_t=\frac{B_{t-h}+B_{t+h}}{2}+2^{-(n+1)/2}Z_t,\qquad h=2^{-n}.
\]

\begin{lemma}[dyadic increments]\label{lem:levylevel}
The increments of this process over the dyadic intervals of level $n$ are i.i.d.\ $N(0,2^{-n})$.
\end{lemma}
\begin{proof}
An interval of variance $2^{-n+1}$ splits into $\Delta/2\pm 2^{-(n+1)/2}Z$. Each half has variance $2^{-n}$ and the two halves are uncorrelated. Induction and Theorem~\ref{thm:gaussvec} give independence.
\end{proof}

\begin{theorem}[existence]\label{thm:levy}
The piecewise-linear interpolations converge uniformly, almost surely, to a Brownian motion on $[0,1]$.
\end{theorem}
\begin{proof}
Level $n$ has at most $2^n$ new displacements, each $N(0,2^{-(n+1)})$. The threshold $a_n=\sqrt{2n}\,2^{-n/2}$ satisfies $a_n^2/(2\sigma^2)=2n$. Proposition~\ref{prop:gausstail} and a union bound produce a probability $O(2^n e^{-2n})$, whose series converges. Theorem~\ref{thm:bc}(a) says that almost surely $a_n$ eventually dominates every displacement. Since $\sum\sqrt n\,2^{-n/2}<\infty$, the interpolations are uniformly Cauchy. The limit is continuous, agrees with the dyadic Gaussians, and therefore has the covariance of Proposition~\ref{prop:bmgauss}.
\end{proof}

\section{Scaling, variation, reflection}

\begin{proposition}[scaling]\label{prop:scaling}
If $B$ is a Brownian motion and $c>0$, then $c^{-1/2}B_{ct}$ is a Brownian motion, and so is $t B_{1/t}$ for $t>0$ with value $0$ at $0$.
\end{proposition}
\begin{proof}
Variances scale by the factor in front of $B$. The inverted process has covariance $st/\max(s,t)=\min(s,t)$. Continuity at $0$ is the same Borel--Cantelli estimate as in Theorem~\ref{thm:levy}.
\end{proof}

\begin{lemma}[fourth moment]\label{lem:squares}
A normal $N(0,\sigma^2)$ has $\mathbb E[\xi^4]=3\sigma^4$, so $\operatorname{Var}(\xi^2)=2\sigma^4$.
\end{lemma}
\begin{proof}
Differentiate $\int e^{-a x^2}\,\dd x=\sqrt{\pi/a}$ twice in $a$, or integrate by parts against the density. Then subtract $(\mathbb E[\xi^2])^2=\sigma^4$.
\end{proof}

\begin{theorem}[quadratic variation]\label{thm:qv}
The sum of squares of the dyadic increments of $B$ on $[0,t]$ converges to $t$ in $L^2$ and almost surely.
\end{theorem}
\begin{proof}
The sum of the means is $t$. The sum of the variances of the squares is $O(2^{-n})$ by Lemma~\ref{lem:squares}. Convergence in $L^2$ follows, and $\sum_n$ of those variances converges, so Theorem~\ref{thm:bc}(a) gives the almost-sure statement.
\end{proof}

\begin{corollary}[variation]\label{cor:variation}
Almost every path has infinite variation and finite quadratic variation on $[0,1]$.
\end{corollary}
\begin{proof}
Finite variation would force the sum of squares of dyadic increments to tend to $0$.
\end{proof}

\begin{theorem}[nowhere differentiable]\label{thm:pwz}
Almost every path is nowhere differentiable on $(0,1)$.
\end{theorem}
\begin{proof}
A derivative at a point would make the dyadic increments in a neighbourhood $O(2^{-n})$, and the sum of their squares would tend to $0$, contradicting Theorem~\ref{thm:qv}.
\end{proof}

\begin{proposition}[martingale property]\label{prop:bmmart}
For $s<t$, $\mathbb E[B_t\mid B_u,\,u\le s]=B_s$.
\end{proposition}
\begin{proof}
The increment $B_t-B_s$ is independent of the past and has mean $0$.
\end{proof}

\begin{proposition}[restart at a level]\label{prop:restart}
Let $\tau_a$ be the first time $B$ hits $a>0$. On $\{\tau_a<\infty\}$ the process $B_{\tau_a+u}-a$ is a Brownian motion independent of the history up to $\tau_a$.
\end{proposition}
\begin{proof}
The dyadic construction after the first dyadic time beyond $\tau_a$ uses fresh normals. Passing to the limit, which Theorem~\ref{thm:levy} justifies uniformly, restarts the increments. The overshoot of the dyadic grid tends to $0$, and the Brownian increment over that vanishing gap tends to $0$.
\end{proof}

\begin{theorem}[reflection]\label{thm:reflection}
For $a>0$ and $t>0$, $P(\max_{s\le t}B_s\ge a)=2P(B_t\ge a)$.
\end{theorem}
\begin{proof}
After the first hit of $a$ the future increments are a fresh Brownian motion, by restarting the dyadic construction. Reflecting that future shows
\[
 P(\tau_a\le t,\,B_t<a)=P(\tau_a\le t,\,B_t>a)=P(B_t>a),
\]
the last equality because a continuous path that ends above $a$ has hit $a$. Adding the two equal pieces gives the result. There is no atom at $a$.
\end{proof}

For $f$ of class $C^2$, Taylor expansion along the dyadic partition and Theorem~\ref{thm:qv} give
\[
 f(B_t)=f(0)+\sum f'(B_{t_{k-1}})\Delta B_k+\tfrac12\int_0^t f''(B_s)\,\dd s
\]
in the limit. The sum is the dyadic stochastic integral. Volume~12 extends it beyond a fixed partition.

\section{Exercises}

\begin{exercise}[$\star\star$]\label{ex:maxjoint}
For $a>0$ and $b\le a$, show $P(M_1\ge a,\,B_1\le b)=P(B_1\ge 2a-b)$.
\end{exercise}

\begin{exercise}[$\star$]\label{ex:bridge}
Let $X_t=B_t-tB_1$. Show that $\operatorname{Cov}(X_s,X_t)=\min(s,t)-st$ and that $X$ is independent of $B_1$.
\end{exercise}

\begin{exercise}[$\star\star$]\label{ex:zeros}
Let $\mathcal Z=\{t\in[0,1]:B_t=0\}$. Show that $\mathcal Z$ is closed, has Lebesgue measure $0$ almost surely, and that $0$ is not isolated in $\mathcal Z$. Deduce that $\mathcal Z$ is uncountable.
\end{exercise}

\begin{exercise}[$\star$]\label{ex:hittingdensity}
Show that the hitting time $T_a$ of level $a>0$ has the law of $a^2/B_1^2$.
\end{exercise}

\begin{exercise}[$\star$]\label{ex:pvariation}
Show that the sum of $p$th powers of dyadic increments on $[0,1]$ tends to $0$ almost surely for $p>2$ and to infinity for $p<2$.
\end{exercise}

\begin{exercise}[$\star$]\label{ex:putcall}
Show that if a call and a put of the same strike $K$ and maturity $T$ are priced so that $C-P\ne S_0-Ke^{-rT}$, the four positions call, put, stock, and cash admit an arbitrage.
\end{exercise}

\begin{exercise}[$\star$]\label{ex:onestep}
Check $\mathbb E[B_t^2]=t$.
\end{exercise}

\section{Hints}

\begin{hint}{ex:maxjoint}
Reflect the path after it first hits $a$.
\end{hint}
\begin{hint}{ex:bridge}
Expand the covariance. The covariance with $B_1$ is zero.
\end{hint}
\begin{hint}{ex:zeros}
$\int_0^1 P(|B_t|<\varepsilon)\,\dd t\to0$. A path that is both positive and negative has a zero between those times. A perfect compact set is uncountable.
\end{hint}
\begin{hint}{ex:hittingdensity}
$\{T_a\le t\}=\{M_t\ge a\}$, and scaling gives $T_a=a^2 T_1$ in law.
\end{hint}
\begin{hint}{ex:pvariation}
Compare expectations with $2^{n(1-p/2)}$.
\end{hint}
\begin{hint}{ex:putcall}
Long call and short put pays $S_T-K$.
\end{hint}
\begin{hint}{ex:onestep}
The variance of $B_t$ is $t$.
\end{hint}

\section{Solutions}

\begin{solution}{ex:maxjoint}
\writeup Reflecting the path after $\tau_a$ sends the endpoint $B_1$ to $2a-B_1$ and preserves the law of the future. The event $\{M_1\ge a,\,B_1\le b\}$ is mapped onto $\{B_1\ge 2a-b\}$.
\end{solution}

\begin{solution}{ex:bridge}
\writeup If $s\le t$, then $\operatorname{Cov}(B_s-sB_1,B_t-tB_1)=\min(s,t)-st$. The covariance of $X_s$ with $B_1$ is $s-s=0$, so Theorem~\ref{thm:gaussvec} gives independence for the Gaussian vector $(X_s,B_1)$.
\end{solution}

\begin{solution}{ex:zeros}
\writeup Continuity makes $\mathcal Z$ closed. The expected Lebesgue measure of $\{|B_t|<\varepsilon\}$ equals $\int_0^1 P(|B_t|<\varepsilon)\,\dd t$, which is the integral of a bounded continuous function of $t$ against the normal law and tends to $0$ with $\varepsilon$ by the bound $P(|B_t|<\varepsilon)\le C\varepsilon t^{-1/2}$ and Theorem~\ref{thm:dct} on $[\delta,1]$, together with a direct estimate near $0$. Therefore $\mathbb E[\lambda(\mathcal Z)]=0$. This uses the integral of Chapter~5, not an external Tonelli theorem. Next, $P(M_\varepsilon>0)=P(\min_{s\le\varepsilon}B_s<0)=1$ by Theorem~\ref{thm:reflection} and symmetry. A continuous path that attains both signs has a zero strictly between those times, so $0$ is not isolated. A nonempty closed set with no isolated point contains a binary tree of distinct points, one for each infinite binary sequence, and is therefore uncountable.
\end{solution}

\begin{solution}{ex:hittingdensity}
\writeup Theorem~\ref{thm:reflection} says $P(T_a\le t)=2P(B_t\ge a)$. Scaling gives $T_a$ the law of $a^2 T_1$. And $P(T_1\le t)=2P(B_1\ge t^{-1/2})=P(B_1^2\ge 1/t)=P(1/B_1^2\le t)$.
\end{solution}

\begin{solution}{ex:pvariation}
\writeup The expected sum of $p$th powers is of order $2^{n(1-p/2)}$. It tends to $0$ for $p>2$ and to infinity for $p<2$. The relative variance tends to $0$, so the almost-sure statement follows from Theorem~\ref{thm:bc} as in Theorem~\ref{thm:qv}.
\end{solution}

\begin{solution}{ex:putcall}
\writeup A long call and a short put pay $S_T-K$. One share financed by borrowing $Ke^{-rT}$ pays the same amount. If the initial prices differ, buy the cheaper side and sell the dearer side.
\end{solution}

\begin{solution}{ex:onestep}
\writeup $B_t$ is centred with variance $t$.
\end{solution}

\section*{Sources and further reading}

\begin{enumerate}
\item N.~Wiener, Differential space, \emph{J.\ Math.\ and Phys.} \textbf{2} (1923), 131--174.
\item P.~L\'evy, \emph{Processus stochastiques et mouvement brownien}, Gauthier-Villars, 1948. The interpolation of Theorem~\ref{thm:levy}.
\end{enumerate}
